% ==============================================================================
% Fichier : front/01_avant_propos.tex
% Titre   : Avant-propos - TPIN
% ==============================================================================

\chapter*{Avant-propos}
\addcontentsline{toc}{chapter}{Avant-propos}
\markboth{Avant-propos}{Avant-propos}

\epigraph{%
    « La technique n'est jamais seulement technique.
    Elle redéfinit l'humain et son rapport au monde. »%
}{Bernard Stiegler}

\vspace{0.5ex}

Ce cours est né d'une double conviction : que l'\termecle{investigation numérique}
constitue aujourd'hui l'une des disciplines les plus stratégiques pour le continent
africain, et que l'Afrique n'a pas à importer une vision de cette discipline, elle
doit en construire une qui lui soit propre, ancrée dans ses réalités juridiques, ses
enjeux économiques et ses capacités institutionnelles.

\bigskip

\textbf{Ce que ce cours n'est pas.} Il ne s'agit pas d'un catalogue d'outils, ni
d'une traduction d'un syllabus occidental. Les ouvrages de référence en langue
anglaise (Casey, Carrier, Carrier, Ligh) font un travail remarquable pour leurs
contextes. Ce cours assume pleinement un ancrage africain et, plus précisément,
camerounais : la Loi~2010/012, le rôle de l'\sigle{ANTIC}, la jurisprudence
naissante des tribunaux, les fraudes \sigle{MoMo},
le \sigle{SIM swapping} en zone \sigle{CEMAC}, voilà les matières premières
de l'investigateur numérique camerounais.

\bigskip

\textbf{Ce que ce cours est.} Le second volume, en trois parties (IV à VI), d'un
traité qui en compte sept (le premier volume est l'UE SEC4031). Il progresse de la
philosophie vers la pratique, de la règle de droit vers l'outil technique, du cas
simple vers la preuve post-quantique. Chaque chapitre comporte des objectifs
mesurables, des exercices gradués, des cas opérationnels ancrés dans la réalité
locale, et une analyse systématique à travers le prisme du \termecle{Trilemme CRO}, cadre théorique original que les travaux de notre équipe au \sigle{LIMSI} ont
formalisé en 2025.

\bigskip

\textbf{Le fil rouge : le Trilemme CRO.} Tout au long de cet ouvrage, vous
rencontrerez l'encadré $\Gamma_{CRO}$ qui interroge systématiquement chaque
technique, chaque protocole, chaque artefact sous trois dimensions :
\begin{itemize}
    \item \textbf{Confidentialité} ($\mathcal{C}$) : la technique protège-t-elle
          les données sensibles des parties impliquées ?
    \item \textbf{Fiabilité} ($\mathcal{R}$) : le résultat est-il intègre,
          reproductible, non altéré ?
    \item \textbf{Opposabilité juridique} ($\mathcal{O}$) : la preuve produite
          sera-t-elle recevable devant un tribunal ?
\end{itemize}
Le Trilemme établit qu'on ne peut pas maximiser simultanément ces trois dimensions.
Comprendre ce compromis, c'est comprendre l'essence même de l'investigation
numérique.

\textbf{Comment utiliser ce cours.} Chaque chapitre s'ouvre sur une \emph{boîte
objectifs} précisant les compétences attendues, et se ferme sur une \emph{boîte
résumé}. Les exercices sont gradués en trois niveaux : \niveaubase\ fondamental,
\niveaumoyen\ intermédiaire, \niveauexpert\ expert. Les codes sources sont testés
et fonctionnels sous Linux (Ubuntu 22.04+) et Windows 10/11.

\bigskip

\begin{flushright}
    \textit{\auteurcours}\\
    \textit{Yaoundé, juin 2026}
\end{flushright}
